\documentclass[11pt]{article}

\usepackage[a4paper,margin=1in]{geometry}
\usepackage{hyperref}
\usepackage{enumitem}
\usepackage{titlesec}

\titleformat{\section}{\large\bfseries}{\thesection}{1em}{}
\setlist[itemize]{noitemsep, topsep=4pt}
\setlist[enumerate]{noitemsep, topsep=4pt}

\title{KCPC-CS111 \\ Supplementary Reading \\ Week 1: Introduction to Competitive Programming, Invariants, and Logic}
\author{King's Competitive Programming Club}
\date{}

\begin{document}

\maketitle

\noindent
This document contains curated supplementary readings for Workshop 1. These texts expand upon the core themes introduced in the session: competitive programming fundamentals, invariant reasoning, epistemic logic (common knowledge), and impartial combinatorial games.

\vspace{1em}
\hrule
\vspace{1em}

\section*{I. Competitive Programming Fundamentals}

\subsection*{Competitive Programmer's Handbook (Antti Laaksonen)}
\url{https://cses.fi/book/book.pdf}

\begin{itemize}
    \item \textbf{Chapter 1: Introduction} Contest environment, basic language features, and input/output optimisation.
    \item \textbf{Chapter 2: Time Complexity} Asymptotic notation ($O$-notation), estimating operation counts per second ($10^8$ ops/sec rule), and detecting TLE bounds.
\end{itemize}

\subsection*{Competitive Programming 4 (Steven Halim, Felix Halim, Suhendry)}
\url{https://cpbook.net/}

\begin{itemize}
    \item \textbf{Chapter 1: Introduction} Contest categories (ICPC, IOI, Codeforces, AtCoder), standard submission verdicts, and typical problem pitfalls.
\end{itemize}

\vspace{1em}
\hrule
\vspace{1em}

\section*{II. Invariants and Mathematical Heuristics}

\subsection*{Problem-Solving Strategies (Arthur Engel)}
\url{https://link.springer.com/book/10.1007/978-0-387-21607-2}

\begin{itemize}
    \item \textbf{Chapter 1: The Invariance Principle} Foundational treatment of invariants, monovariants, and coloring arguments in discrete dynamical processes.
\end{itemize}

\subsection*{The Art and Craft of Problem Solving (Paul Zeitz)}
John Wiley \& Sons, 2006.

\begin{itemize}
    \item \textbf{Chapter 3: Tactics for Solving Problems} Symmetry, parity, extremal thinking, and simplifying assumptions.
\end{itemize}

\vspace{1em}
\hrule
\vspace{1em}

\section*{III. Epistemic Logic and Common Knowledge}

\subsection*{Reasoning About Knowledge (Fagin, Halpern, Moses, Vardi)}
MIT Press, 2003. \url{https://mitpress.mit.edu/9780262562003/reasoning-about-knowledge/}

\begin{itemize}
    \item \textbf{Chapter 1: Introduction} The distinction between mutual knowledge and common knowledge; the formal mud-on-the-forehead puzzle.
    \item \textbf{Chapter 2: A Model for Knowledge} Kripke structures and epistemic possible worlds.
\end{itemize}

\vspace{1em}
\hrule
\vspace{1em}

\section*{IV. Combinatorial Game Theory}

\subsection*{Game Theory (Thomas S. Ferguson)}
\url{https://www.math.ucla.edu/~tom/Game_Theory/comb.pdf}

\begin{itemize}
    \item \textbf{Part I: Impartial Combinatorial Games} Subtraction games, normal-play conventions, $N$- and $P$-positions, and backward induction on directed acyclic graphs.
\end{itemize}

\subsection*{CP-Algorithms: Games on Arbitrary Graphs}
\url{https://cp-algorithms.com/game_theory/games_on_graphs.html}

\begin{itemize}
    \item Topological determination of winning and losing states via retro-analysis / backward induction.
\end{itemize}

\end{document}
